%%%
% src::
%     urls = https://tex.stackexchange.com/a/654470,
%            https://tex.stackexchange.com/a/761377/6880
%%%

\documentclass[border = 3pt]{standalone}

%%%
% Le \pack ''nicematrix'' est le couteau suisse des matrices.
%%%
\usepackage{nicematrix}

%%%
% Le \pack ''tikz'' va permet de décorer le tableau d'adjacence.
%%%
\usepackage{tikz}

\begin{document}

\renewcommand{\arraystretch}{1.2}

\begin{NiceTabular}{>{\bfseries}*{6}{c}}[
  hvlines,
  corners = NW,
  name    = adj-matrix
]
  \RowStyle{\bfseries}
    & A & B & C & D & E  \\
  A &   &   & 1 & 1 &    \\
  B &   &   &   & 1 & 1  \\
  C & 1 &   &   &   & 1  \\
  D & 1 & 1 &   &   &    \\
  E &   & 1 & 1 &   &
\end{NiceTabular}

%%%
% Le point d'ancrage ''shipout/background'' du noyau \latex exécute
% du code au moment de l'envoi de la page vers le PDF, en plaçant
% les éléments sur une couche située derrière tout le contenu déjà
% présent.
%%%
\AddToHook{shipout/background}{
  \begin{tikzpicture}[remember picture, overlay]
    \draw[red!60, ->, shorten <= 1pt, shorten >= 1pt]
      (adj-matrix-4-1.east)
      -| node[pos=0.5, name=corner] {}
      (adj-matrix-1-2.south);
    \fill[red!10] (corner) circle (5pt);
  \end{tikzpicture}
}

\end{document}
